% Avsnitt 11.1-11.6, ur lectures/L11/appendix/a_architecture_and_register_map.md.

\section{Kontrollerns arkitektur}
\label{c11:sec:arch}

Varje regel från \chapref{c10:ch} avbildas på ett av fem block:
\begin{itemize}
  \item \textbf{\code{bit_timer}} äger ingenting om \emph{vad} som sänds, bara \emph{när}. Den är det
    enda i den här konstruktionen som ens gränsar till en klockdomänövergång: den frilöper
    kontinuerligt så snart den är aktiverad, och resynkroniserar exakt en gång per ram, vid SOF
    (\secref{c10:sec:stuffing}:s resonemang om klockåtervinning, tillämpat).
  \item \textbf{\code{tx_shift_reg}} och \textbf{\code{rx_shift_reg}} äger bitstoppningen och
    ingenting annat: de serialiserar eller deserialiserar upp till 8 riktiga bitar i taget och sätter
    transparent in eller tar bort stoppbitar. De har ingen aning om huruvida bitarna de skiftar är
    ett ID, en databyte eller ett CRC-värde; det är \code{can_controller}s jobb.
  \item \textbf{\code{crc15}} äger CRC-15-beräkningen, en bit i taget, utan någon föreställning om
    ”ramar” eller ”fält” den heller; den ackumulerar bara den bit den får, när den blir tillsagd.
  \item \textbf{\code{can_controller}} är det enda block som känner till det faktiska CAN-ramformatet:
    det sekvenserar genom SOF, arbitrering, kontrollfältet, data, CRC, ACK och EOF, i ordning, driver
    de andra fyra blocken och läser deras utgångar. Detektering av förlorad arbitrering bor också
    här, eftersom den specifikt handlar om att jämföra vad den här kontrollern bestämde sig för att
    sända mot vad bussen faktiskt visar.
\end{itemize}

\lecfig[0.92]{lectures/L11/appendix/images/can_controller_diagram.png}{\code{can_controller}s
blockarkitektur.}{c11:fig:arch}

Två konventioner håller bilden läsbar. \code{clock} når varje delblock, och den synkroniserade
\code{reset_s2_n} når alla utom de två \code{meta_prev}-instanser som producerar den och dess
bussmotsvarighet; båda lämnas oritade, som brukligt i ett blockschema. Varje utgångsport
(\code{tx_done}, \code{tx_bus}, \code{bus_en}, \code{rx_id}, \code{rx_dlc}, \code{rx_data},
\code{rx_valid}, \code{error}) drivs av ramtillståndsmaskinen och av ingenting annat.

Tre av de elva interna kopplingarna är värda att läsa två gånger, eftersom det är de som folk ritar
om fel ur minnet:
\begin{itemize}
  \item De två linjerna vid \code{crc15} går i \textbf{motsatta} riktningar. \code{clear},
    \code{enable} och \code{data} går upp in i den, en riktig bit i taget; \code{crc} och
    \code{valid} kommer tillbaka ner ur den. (\code{clear} är \chapref{c13:ch}:s reset per ram
    \textemdash{} motorn måste börja varje ram på noll, och en avbruten ram skulle annars lämna
    rester efter sig; \chapref{c16:ch} driver den från \code{STATE_START}.)
  \item \code{rx_shift_reg} har \textbf{tre} linjer som anländer: \code{sample} från
    \code{bit_timer}, \code{rx_bus_s2} från synkroniseraren som \chapref{c12:ch} lägger till, och
    \code{bit_count} plus \code{enable} från ramtillståndsmaskinen. Utan de två första har den ingen
    datakälla alls.
  \item \code{tx_shift_reg} och \code{rx_shift_reg} kopplas aldrig till varandra, och ingen av dem
    kopplas till \code{crc15}. Ramtillståndsmaskinen bär varje bit mellan dem. \code{rx_shift_reg}
    är på samma sätt det enda delblock som över huvud taget läser bussen, och till och med det läser
    den synkroniserade kopian snarare än porten; se \secref{c11:sec:sync}.
\end{itemize}

\paragraph{Hur de nästlas}
Diagrammet läses som en lagerbild, med det billigaste bekymret överst, men den faktiska
VHDL-hierarkin är inneslutning, inte en pipeline: \code{can_controller} instansierar alla fyra av
\code{bit_timer}, \code{crc15}, \code{tx_shift_reg} och \code{rx_shift_reg} inuti sig själv, vid
sidan av den \code{meta_prev}-synkroniserare som \chapref{c12:ch} lägger till.
\code{can_controller} är därför den enda entitet du någonsin instansierar för hand, och hela
kontrollern är en komponent utifrån sett, exakt som registerbanken i \chapref{c18:ch} kommer att se
den.

Två lager till sitter sedan ovanför den: \code{register_bank} (\chapref{c18:ch}), som gör om de
portarna på registernivå till den registerkarta en drivrutin pollar, och \code{can_spi_node}
(\chapref{c19:ch}), toppnivån som lägger till SPI-transporten och kortets pinnar. Båda hör till att
nå kontrollern, inte till att tala CAN, vilket är varför de dyker upp i de två sista kapitlen snarare
än här.

Den här uppdelningen \textemdash{} ett block per \emph{ansvar}, inte ett block per \emph{ramfält}
\textemdash{} är medveten: \code{tx_shift_reg} återanvänds för ID:t, kontrollfältet, varje databyte
och CRC-värdet, helt enkelt omladdat med olika data och olika bitantal varje gång, i stället för att
skriva nästan identisk skiftlogik fem separata gånger.

% ------------------------------------------------------------------------------------------
\section{Att följa en ram genom blocken}
\label{c11:sec:trace}

Att namnge de fem blocken är inte samma sak som att veta hur de samarbetar. Det här avsnittet tar en
konkret ram, ID \code{0x123}, DLC \code{2}, databytesen \code{A1 B2}, och följer den genom
arkitekturen ovan. Det finns ingen VHDL här; frågan är bara vilket block som gör vad, och vad som
passerar mellan dem.

\paragraph{Först, det enda block som räknar tid}
\code{bit_timer} delar upp 50 MHz-klockan i bitperioder och avger två pulser per period:
\code{sample}, 70~\% in i perioden, när nivån på bussen hunnit stabilisera sig, och \code{bit_done},
allra sist, när den aktuella biten är klar. Varje annat block agerar på en av de två pulserna, och
inget av dem räknar klockcykler mot bussen själv. Sändning rör sig på \code{bit_done}, vid gränsen
där en ny bit börjar. Mottagning rör sig på \code{sample}, mitt i en bit som redan ligger där. Just
den skillnaden är varför en mottagare läser det värde en sändare lagt ut, i stället för att fånga
det mitt i en förändring.

\paragraph{Sändning: 60 bitar, 50 av dem genom ett register}
Med fältbredderna från \chapref{c10:ch} är den här ramen $1$ (SOF) $+\ 11$ (ID) $+\ 1$ (RTR)
$+\ 2$ (IDE, r0) $+\ 4$ (DLC) $+\ 16$ (data) $+\ 15$ (CRC) $=$ \textbf{50 bitar i det stoppade
området}, följt av en ostoppad svans på $1 + 1 + 1 + 7 =$ \textbf{10 bitar}. Sextio bitar innan någon
stoppbit satts in.

De 50 bitarna når bussen som \textbf{åtta laddningar av ett enda \code{tx_shift_reg}}:

\begin{narrowtable}{clr}
  \thead{\#} & \thead{Grupp} & \thead{Bitar} \\
  \midrule
  1 & SOF & 1 \\
  2 & Identifieraren, hög del & 8 \\
  3 & Identifieraren, låg del + RTR & 4 \\
  4 & IDE, r0, DLC & 6 \\
  5 & Databyte 0 (\code{A1}) & 8 \\
  6 & Databyte 1 (\code{B2}) & 8 \\
  7 & CRC, hög del & 8 \\
  8 & CRC, låg del & 7 \\
\end{narrowtable}

\code{can_controller} laddar om samma 8-bitarsregister åtta gånger, varje gång med ett annat värde
och ett annat bitantal. Det är påståendet ”ett block per ansvar” ovan, gjort konkret: identifieraren
är inget fält för skiftregistret, den är grupp 2. De återstående 10 svansbitarna når aldrig något
skiftregister alls; de är fasta recessiva nivåer, och \code{can_controller} driver dem direkt medan
den räknar \code{bit_done}-pulser.

Inom en grupp händer fyra saker per bit:
\begin{itemize}
  \item \code{tx_shift_reg} producerar nästa bit. Laddningen presenterar den första omedelbart; varje
    \code{bit_done} därefter skiftar ut nästa. Registret avgör också, på egen hand, om en stoppbit
    ska sättas in, och annonserar på \code{stuff} vilken sorts bit det här var.
  \item \code{can_controller} gör om biten till bussnivåer. Dominant betyder \code{bus_en = '1'} med
    \code{tx_bus = '0'}; recessiv betyder \code{bus_en = '0'}, vilket släpper bussen. En nod driver
    aldrig recessivt aktivt (\secref{c10:sec:can}:s wired-AND).
  \item \code{crc15} matas med samma bit, men bara när den var en riktig. Stoppbitar hoppas över.
    \code{crc15} har ingen aning om vad en stoppbit är; att bli tillsagd när den \emph{inte} ska
    ackumulera är hela dess gränssnitt på den sidan.
  \item \code{can_controller} jämför \code{rx_bus} mot biten den just drev. Den jämförelsen är
    arbitreringen, och den stannar här snarare än i ett delblock eftersom den behöver både vad den
    här noden avsåg och vad bussen faktiskt visar, och \code{can_controller} är det enda block som
    håller båda.
\end{itemize}

Efter grupp 6 håller \code{crc15} CRC:n över allt som sänts hittills. \code{can_controller} läser
den och matar tillbaka den rakt in i \code{tx_shift_reg} som grupp 7 och 8. Lägg märke till vad det
betyder: motorn som beräknar CRC:n och registret som sänder den kopplas aldrig till varandra.
\code{can_controller} bär värdet över, vilket är varför inget av blocken behöver veta att ett
CRC-fält finns.

\paragraph{Mottagning: samma väg, speglad}
\code{rx_shift_reg} samplar \code{rx_bus} på \code{bit_timer.sample}, tar bort stoppbitar och
rapporterar varje riktig bit på \code{real_bit} och \code{real_bit_valid}. De två signalerna matar
\code{crc15} precis som \code{tx_shift_reg}s gör vid sändning: samma instans, samma regel om att
bara ackumulera riktiga bitar, ingen lägesomkoppling någonstans (vilket är \chapref{c13:ch}:s poäng
om att en motor gör båda jobben). Fullbordade grupper anländer på \code{data}, och
\code{can_controller} skivar tillbaka dem till \code{rx_id}, \code{rx_dlc} och \code{rx_data}. En
stoppbit som uteblir höjer \code{stuff_error}, och \code{rx_shift_reg} stannar där: den rapporterar
brottet utan att avgöra vad det betyder. Att avgöra är \code{can_controller}s jobb, och den bestämmer
sig för att avbryta ramen och höja \code{error}.

\paragraph{Vad genomgången visar}
Inget delblock vet vilket fält det hanterar. Vart och ett rapporterar fakta om \emph{bitar}: en är
klar, den där sattes in, en stoppbit saknades, den här gruppen är fullbordad. \code{can_controller}
levererar varenda del av meningen. Tre saker följer, och alla tre är värda att ha:
\begin{itemize}
  \item De fyra delblocken går att testa utan CAN. Deras testbänkar i \chapref{c12:ch} till
    \ref{c15:ch} driver bitmönster och kontrollerar bitmönster; inte en enda av dem behöver en hel
    ram för att vara meningsfull.
  \item \code{can_controller} är det enda block som skulle ändras om ramformatet gjorde det. Utökade
    29-bitars identifierare skulle till exempel betyda andra grupper, inte andra block.
  \item Varje koppling mellan block är en signal på en bit eller en byte med exakt en skrivare. Det
    finns ingen delad buss eller gemensam struktur som två block båda griper in i, vilket är det som
    håller kopplingsarbetet i \chapref{c12:ch} till \ref{c15:ch} vid en port map i taget snarare än
    en omkonstruktion varje kapitel.
\end{itemize}

% ------------------------------------------------------------------------------------------
\section{Kontrollerns gränssnitt}
\label{c11:sec:interface}

\code{can_controller}s portar är vad omvärlden ser, och de ligger fast resten av boken:
\chapref{c12:ch} till \ref{c17:ch} lägger till delblock \emph{inuti} den här entiteten utan att ändra
en enda port. Två saker kopplas till den, och varje port hör till en av dem.

\lecfig[0.66]{lectures/L11/appendix/images/can_controller.png}{\code{can_controller}s femton
portar.}{c11:fig:ports}

\paragraph{Bussidan, fem portar}
Det är de portar som når omvärlden som fysiska pinnar; de tre första vetter mot CAN-transceivern:
\begin{itemize}
  \item \code{rx_bus} är den aktuella bussnivån, en bit, direkt från transceivern och asynkron mot
    \code{clock}. Kontrollern synkroniserar den själv, med en \code{meta_prev}-instans som
    \chapref{c12:ch} bygger, så att kortets toppnivå kan koppla den här porten till en pinne och inte
    tänka mer på det. Kontrollern läser den för att ta emot, och också för att övervaka sina egna
    sändningar under arbitreringen.
  \item \code{tx_bus} är biten den här noden vill driva, och \code{bus_en} säger om den driver över
    huvud taget. Paret modellerar beteendet hos öppen dränering: \code{bus_en = '1'} med
    \code{tx_bus = '0'} drar bussen dominant, och \code{bus_en = '0'} släpper den så att en annan nod
    kan göra det. En nod driver aldrig recessivt aktivt; den släpper helt enkelt taget. Det är det
    som får \secref{c10:sec:can}:s wired-AND att fungera.
  \item \code{clock} är 50 MHz-systemklockan, och \code{reset_n} en \textbf{rå}, aktivt låg reset,
    direkt från en pinne. Precis som \code{rx_bus} är den asynkron, och precis som \code{rx_bus}
    synkroniserar kontrollern den själv. Den synkroniserade kopian är en intern signal som heter
    \code{reset_s2_n}, namnet övningen \code{reset_sync} gav den i \chapref{c4:ch} och namnet varje
    delblock i projektet fortfarande tar.
\end{itemize}

\paragraph{Sändningssidan, fem portar}
Det är dem en anropare fyller i för att sända en ram:
\begin{itemize}
  \item \code{tx_id}, \code{tx_dlc} och \code{tx_data} bär själva ramen, typade som \code{id_t},
    \code{dlc_t} och \code{data_t} från \code{can_def}. \code{tx_data} håller byte 0 i sina mest
    signifikanta bitar, i samma ordning som bussen sänder dem.
  \item \code{tx_req} begär sändning. Kontrollern agerar på den bara när den själv är i vila
    \textbf{och bussen också har varit ledig}, eftersom en nod inte får börja driva ovanpå en ram som
    redan pågår (\chapref{c16:ch}). En begäran som höjs medan bussen är upptagen kastas i stället för
    att köas, så anroparen får begära igen.
  \item \code{tx_done} pulsar högt under en cykel så snart den här nodens egen ram har sänts färdigt.
\end{itemize}

\paragraph{Mottagningssidan, fem portar}
De speglar sändningssidan exakt:
\begin{itemize}
  \item \code{rx_id}, \code{rx_dlc} och \code{rx_data} håller den senast mottagna ramen, samma typer
    och samma byteupplägg som sina \code{tx_}-motsvarigheter.
  \item \code{rx_valid} pulsar högt under en cykel när en fullständig, giltig ram har anlänt.
  \item \code{error} är den enda port som \emph{inte} är en puls: den går hög vid förlorad
    arbitrering, en stoppningsöverträdelse, en misslyckad CRC-kontroll eller en mottagen fjärram (en
    recessiv RTR-bit), och stannar hög tills nästa ram börjar. En anropare som pollar den då och då
    ser ändå felet.
\end{itemize}

Fyra saker med det här gränssnittet är värda att lägga märke till nu, eftersom de formar allt som
följer:
\begin{itemize}
  \item \textbf{Det finns ingen serialisering i det.} Anroparen lämnar över en identifierare, en
    längd och bytes. Varenda del av ramningen, stoppningen och CRC:n som \chapref{c10:ch} beskrev
    sker inuti, på vägen till \code{tx_bus}. Det är hela poängen med att bygga en kontroller.
  \item \textbf{Sändning och mottagning är symmetriska men inte samtidiga.} En nod sänder antingen
    sin egen ram eller tar emot någon annans, aldrig båda, så gränssnittets två halvor är aldrig
    aktiva samtidigt.
  \item \textbf{\code{tx_done}, \code{rx_valid} och \code{error} är hela statusgränssnittet.} Varje
    statusbit registerkartan erbjuder en drivrutin (\chapref{c18:ch}) härleds ur någon av de tre.
  \item \textbf{Varje port på registersidan antas synkron mot \code{clock}.} \code{rx_bus} är
    undantaget, och det hanterar kontrollern själv med \code{meta_prev}. Resten, sändningsportarna en
    anropare driver och status- och mottagningsportarna den läser, förutsätter en anropare i samma
    klockdomän, vilket är exakt vad \code{register_bank} (\chapref{c18:ch}) är. En anropare i en
    \emph{annan} domän behöver en handskakning framför dem, inte en bredare \code{meta_prev}: att
    köra 64 bitar \code{tx_data} genom en bit-för-bit-synkroniserare tar bort metastabiliteten men
    låter fortfarande kontrollern låsa ett värde som aldrig skrevs, samma skevhetsproblem som
    \chapref{c19:ch} lägger ut för omkopplarbanken. En synkroniseringskedja är rätt verktyg för en
    bit vars grannar är oberoende, och fel verktyg för ett värde vars bitar måste stämma överens.
\end{itemize}

% ------------------------------------------------------------------------------------------
\section{Att skriva \code{can_controller.vhd}}
\label{c11:sec:writing}

Det här kapitlet skriver filen från grunden, till \file{controller/can_controller.vhd}, bredvid den
\file{can_def.vhd} som skrevs i \chapref{c10:ch}.

\textbf{Entiteten} deklarerar de femton portarna som beskrevs ovan, i exakt den här ordningen, vilket
är det \code{can_controller_tb} binder till:

\begin{textdrawing}
 1 clock       4 tx_id     7 rx_bus     10 bus_en    13 rx_data
 2 reset_n     5 tx_dlc    8 tx_done    11 rx_id     14 rx_valid
 3 tx_req      6 tx_data   9 tx_bus     12 rx_dlc    15 error
\end{textdrawing}

\textbf{Varje ingång först, sedan varje utgång}, vilket är den här bokens konvention överallt och är
varför \code{rx_bus} sitter sjua snarare än bredvid de bussutgångar den begreppsligt hör ihop med.
Läs blockschemat ovan rakt ned för dess vänstra kolumn och sedan dess högra, så får du 1 till 15 i
ordning. Typerna kommer från \code{can_def}, så filen öppnar med \code{use work.can_def.all;} vid
sidan av det vanliga \code{ieee.std_logic_1164}, och portlistan blandar rent \code{std_logic} med
paketets subtyper. De första raderna sätter mönstret för alla femton:

\begin{vhdlcode}
library ieee;
use ieee.std_logic_1164.all;
use work.can_def.all;

entity can_controller is
    port(clock   : in std_logic;
         reset_n : in std_logic;
         tx_req  : in std_logic;
         tx_id   : in id_t;
         ...
\end{vhdlcode}

\begin{warning}
Få \textbf{ordningen} exakt rätt, eftersom det är den allt binder till: \code{can_controller_tb}
associerar sina portar positionellt, och det gör varje instansiering i den här boken. Namnen är
dina. Ett omkastat par analyserar alldeles utmärkt i dag, och till skillnad från ett felstavat namn
kommer det inte att misslyckas med att analysera i \chapref{c17:ch} heller. Det kommer i tysthet att
koppla ihop fel signaler, och testbänken kommer att rapportera beteende som inte går ihop, sex
kapitel senare.
\end{warning}

\textbf{Arkitekturen börjar tom.} Inte en platshållare som ska bytas ut senare, utan genuint tom:

\begin{vhdlcode}
architecture behaviour of can_controller is
begin
end architecture;
\end{vhdlcode}

Det analyserar och elaborerar rent, och \code{make build} bekräftar att det gör det. Skriv det,
kontrollera det, och titta på det ett ögonblick: en entitet utan något beteende alls är fortfarande
en fullständig designenhet. Sedan lägger resten av projektet in det första inuti den.

\paragraph{Varför skriva den nu, tom?}
Därför att gränssnittet är ett beslut, och det är ett det här kapitlet är i läge att fatta. Du vet
vad en CAN-ram innehåller (\chapref{c10:ch}), du vet vilka block som ska göra jobbet
(\secref{c11:sec:arch}), och du vet vad en drivrutin så småningom behöver läsa och skriva
(\secref{c11:sec:regmap}). Allt \chapref{c12:ch} och framåt lägger till är \emph{beteende bakom
portar som redan finns}.

Alternativet, att skriva delblocken först och upptäcka toppnivåns gränssnitt på slutet, är så
konstruktionen skulle växa fram om ingen hade tänkt på den i förväg. Det betyder också att formen på
det du bygger förblir ett mysterium ända till det allra sista kapitlet.

\paragraph{Vad varje kapitel lägger till}
Härifrån bygger varje kapitel ett block och instansierar det här, så att arkitekturen fylls i
inifrån. \Chapref{c12:ch} gör två, eftersom båda är små:

\begin{booktable}{@{}p{6em}L@{}}
  \thead{Kapitel} & \thead{Lägger till i \code{can_controller}} \\
  \midrule
  \ref{c12:ch} & \code{meta_prev} två gånger, som synkroniserar \code{reset_n} och \code{rx_bus},
    sedan \code{bit_timer} \\
  \ref{c13:ch} & \code{crc15}, och signalerna som matar den en bit i taget \\
  \ref{c14:ch} & \code{tx_shift_reg} \\
  \ref{c15:ch} & \code{rx_shift_reg} \\
  \ref{c16:ch} & Den ramsekvenserande tillståndsmaskinen, sändningsvägen \\
  \ref{c17:ch} & Mottagningsvägen och arbitreringen, som fullbordar konstruktionen \\
\end{booktable}

Filen måste analysera rent i vart och ett av de stegen, inte bara på slutet. \code{make build}
kontrollerar varje modul för sig så snart den finns, så en entitet som inte kompilerar fångas i det
kapitel som skrev den.

% ------------------------------------------------------------------------------------------
\section{Vad kontrollern synkroniserar, och varför det är kontrollerns jobb}
\label{c11:sec:sync}

Ett beslut följer direkt ur den sista punkten i \secref{c11:sec:interface}, och det är värt att göra
uttryckligt här eftersom varje senare kapitel hänger på det. \code{rx_bus} anländer från en
transceiver på ett annat chip, och \code{reset_n} från en knapp eller en övervakare, så ingenting i
den här FPGA:n har någonsin styrt någon av dem. Att sampla den ena direkt in i \code{clock}s domän
kan lämna en vippa metastabil, problemet \chapref{c4:ch} löste med två vippor i serie.

De två är kontrollerns \emph{enda} asynkrona ingångar, och de är asynkrona vem som än instansierar
den, så \textbf{kontrollern synkroniserar dem själv} i stället för att lämna det till kortets
toppnivå. Den som integrerar kopplar båda rakt till pinnar och tänker inte mer på det. Allt annat på
entiteten är registergränssnittet, som är ett annat kontrakt: ett som en anropare i den här
klockdomänen uppfyller gratis, och som en anropare i en annan domän måste uppfylla med en
handskakning snarare än en bredare synkroniserare.

Internt heter de synkroniserade kopiorna \code{reset_s2_n} och \code{rx_bus_s2}, enligt
\chapref{c4:ch}:s konvention från \code{reset_sync}, och det är de namn varje delblock från
\chapref{c12:ch} och framåt tar. \textbf{Inget block läser någonsin \code{reset_n} eller
\code{rx_bus} direkt}, arbitreringsjämförelsen i \chapref{c17:ch} inräknad. Modulen som producerar
dem, \code{meta_prev}, är \chapref{c12:ch}:s, vid sidan av bittimern; det här kapitlet slår fast att
kontrollern äger problemet, och nästa bygger det som löser det.

% ------------------------------------------------------------------------------------------
\section{En första titt på registerkartan}
\label{c11:sec:regmap}

En parallellklass skriver en C++-drivrutin mot den här kontrollern, och talar med den genom
\textbf{register som nås över SPI} snarare än via VHDL-portar direkt. Du bygger lagren som
presenterar de registren: \code{register_bank} i \chapref{c18:ch} och \code{spi_reg_bridge} i
\chapref{c19:ch}. Hela kartan och transaktionsformatet finns i bilaga~\ref{app:protocol}, och
tillsammans är de hela kontraktet mellan de två klasserna.

Den är värd en förhandstitt här, sju kapitel i förväg, eftersom det är exakt samma information som
\code{can_controller}s portar redan bär \textemdash{} de senare lagren ändrar informationens
\emph{form}, aldrig dess innehåll:

\begin{booktable}{@{}p{14em}L@{}}
  \thead{Begrepp (\code{can_controller})} & \thead{Register i kartan} \\
  \midrule
  \code{tx_id}, \code{tx_dlc}, \code{tx_data} & \code{TX_ID}, \code{TX_DLC},
    \code{TX_DATA_LO}/\code{HI} \\
  \code{tx_req}, \code{tx_done} & \code{TX_SEND}, en bit i \code{STATUS} \\
  \code{rx_id}, \code{rx_dlc}, \code{rx_data} & \code{RX_ID}, \code{RX_DLC},
    \code{RX_DATA_LO}/\code{HI} \\
  \code{rx_valid} & en bit i \code{STATUS}, nollställd via \code{RX_ACK} \\
  \code{error} & \code{ERROR_FLAGS} \\
\end{booktable}

Ha den här motsvarigheten i minnet genom \chapref{c19:ch}: allt parallellklassens drivrutin läser
eller skriver är bara den här kontrollerns eget gränssnitt på portnivå, inslaget i register.
